\documentclass{article}
\usepackage{graphicx, amsmath, amssymb, mathrsfs, pgfplots, enumerate, tikz, wrapfig, tcolorbox, forest}
\usepackage[a4paper, top=2cm, bottom=3cm, left=1.5cm, right=1.5cm]{geometry}
\usepackage[vlined, boxruled]{algorithm2e}
\usepackage{listings}
\usepackage{xcolor}
\title{Relazione di laboratorio ASD}
\author{Tobia Viola Romanese e Mauro Brancato}
\date{2025-2026}

\begin{document}

\maketitle







\tableofcontents
\clearpage

\section{Obiettivo}
L'obiettivo di questo progetto è quello di implementare tre tipologie di alberi binari di ricerca: \textbf{Adelson-Velsky and Landis Tree (AVL)}, \textbf{Red Black Tree (RBT)}, \textbf{Binary Search Tree (BST)} e misurare i tempi d'esecuzione
dei loro metodi per verificare sperimentalmente la loro complessità teorica in termini asintotici.
\section{Organizzazione del progetto}
Il progetto è suddiviso nel codice in 4 parti:
\begin{itemize}
    \item \textbf{Le strutture dati}: Abbiamo diviso in tre cartelle separate le varie strutture dati con le loro implementazioni.
    \item \textbf{TreeManager.py}: Abbiamo costruito un TreeManager per gestire le varie strutture dati. Questa classe ci permette di rendere il codice più leggibile, ma anche più scalabile introducendo vantaggi di OOP.
    \item \textbf{Main.py}: È il cuore pulsante del progetto, contiene l'inizializzazione delle varie strutture dati, la creazione della serie geometrica per un'analisi più accurata e la gestione dei vari alberi tramite il TreeManager.
    \item \textbf{grafici.py}: È un file che ci permette, grazie alla libreria \verb|matplotlib.pyplot as plt| di creare dei grafici da file \verb|.csv|.
\end{itemize}
\subsection{Strutture Dati}
Nel progetto abbiamo implementato tre strutture dati per gestire gli alberi: \textbf{Adelson-Velsky and Landis Tree (AVL)}, \textbf{Red Black Tree (RBT)}, \textbf{Binary Search Tree (BST)}. La struttura dati portante del progetto è anche quella più semplice, cioè i \textbf{BST}. Le altre due classi che abbiamo implementato, infatti, ereditano le proprietà di quest'ultima; questa scelta ha una motivazione pratica (hanno dei metodi in comune che non hanno bisogno di modifiche), ma anche una motivazione teorica essendo sia \textbf{RBT} e \textbf{AVL} anch'essi alberi binari di ricerca ma con delle condizioni e specifiche ben precise per gestire l'\textbf{altezza}.
\subsection{TreeManager}
All'interno del progetto per introdurre alcuni concetti imparati nel primo semestre abbiamo deciso di implementare una classe chiamata TreeManager per sfruttare i vantaggi della OOP. La classe ha due variabili d'istanza: \verb|keys_array| e \verb|m|. La prima, \verb|keys_array| è un array che contiene tutte le chiavi che da sfruttare per le misurazioni. L'array è diviso in due, nella prima parte ci sono le chiavi che sono inserite nell'albero, nella seconda metà ci sono le chiavi non ancora inserite che possono essere effettuate per la misurazione dell'inserimento. Per dividere le due parti abbiamo bisogno di un indice che scorre man mano che riempiamo l'albero inserendo le chiavi ed è qui che entra in gioco la seconda variabile d'istanza \verb|m|, un intero che separa le chiavi esterne da quelle già inserite nell'albero.
In questa classe abbiamo costruito due metodi principali: \verb|create_tree| e \verb|measure_insert|.\\\\
Il metodo \verb|create_tree| riceve in input un albero (\verb|tree|) e il suo scopo è quello di riempire la struttura dati in input con le chiavi all'interno della porzione di array $key_array[m:n]$ con $n = length(keys_array)$. Queste chiavi vengono selezionate in modo casuale così da garantire una distribuzione casuale uniforme. Una volta selezionata la chiave, viene creato il nodo e viene inserito nell'albero. Ora la chiave selezionata viene spostata all'interno della porzione dell'array $key_array[0:m-1]$ perché ora la chiave fa parte dell'albero e di conseguenza aumentiamo anche l'indice \verb|m| di uno.\\\\
Il metodo \verb|measure_insert| riceve in input un albero (\verb|tree|) e un numero di iterazioni (\verb|batch_size|) che corrisponde al numero di tempi raccolti per una certo albero di dimensione n. All'inizio del metodo viene creato il vettore che conterrà tutte i tempi misurati, successivamente c'è il ciclo for che costituisce il corpo del metodo. All'interno del ciclo viene scelta la chiave da inserire e viene creato il nodo corrsipondente da inserire nell'albero, viene fatto partire il timer, si effettua l'inserimento e il timer viene fermato. Il tempo calcolato viene inserito nel vettore dei tempi e il vettore \verb|keys_array| viene sistemato a dovere scambiando la chiave appena inserita nell'albero nella prima parte dell'array. Sempre all'interno del ciclo per ottenere delle misurazioni corrette per una dimensione fissa n dobbiamo rimuovere una chiave in modo casuale dall'albero. 





\end{document}